% ============================================================================
% Chapter 2, Section 3: Related Work
% ============================================================================

\subsection{Related Work}\label{sec:related_work}

We survey five lines of research that collectively inform the design of our benchmark, harness, and AutoEvolve algorithm: embodied agents, experiential learning, recursive language models, self-adaptation methods, and the emerging field of meta-harnesses.

\subsubsection{Embodied Agents}\label{sec:embodied_agents}

The deployment of language models as controllers for embodied agents has produced a rich body of work exploring how latent world knowledge can be translated into effective physical or virtual action. Voyager~\citep{wang2023voyager} demonstrated open-ended skill acquisition in Minecraft through an automatic curriculum and a growing code library, achieving lifelong learning without manual reward engineering. However, Voyager relies on a programmatic API that provides structured state information (inventory contents, nearby blocks, entity positions) rather than raw perceptual input, sidestepping the perception challenges central to RPG gameplay.

SayCan~\citep{brohan2023can} grounded language model instructions in robotic affordances, showing that LLMs can select contextually appropriate actions when paired with learned value functions---but operates in short-horizon manipulation tasks rather than extended autonomous episodes. Code as Policies~\citep{liang2023code} generalized this approach by having language models generate executable robot control code, demonstrating compositional generalization to novel instructions. GITM (Ghost in the Minecraft)~\citep{zhu2023ghost} combined LLMs with text-based knowledge and memory to create generally capable open-world agents, introducing a knowledge graph for tracking world state. Steve-Eye~\citep{zheng2024steveeye} extended this line of work by equipping LLM-based agents with visual perception in Minecraft, moving closer to the perceptual challenges we face in Pok\'{e}mon.

Generative Agents~\citep{park2023generative} demonstrated that LLMs can produce believable human behavior in a simulated town environment when equipped with memory retrieval, reflection, and planning mechanisms. Their memory architecture---combining recency, importance, and relevance-based retrieval---has influenced subsequent work on agent memory systems, including our own hierarchical memory design in PokeAgent.

A common thread across these works is that none operates from purely perceptual input (raw frames) in a long-horizon RPG setting requiring sustained reasoning over many hours. The closest analogue to our setting is Voyager, but it benefits from a structured API, operates in episodes of minutes rather than hours, and does not contend with the partially observable, narratively structured progression that characterizes Pok\'{e}mon RPGs. Our work bridges this gap by developing a harness that processes raw visual observations alongside structured state text and supports multi-hour continuous operation.

\subsubsection{Experiential Learning}\label{sec:experiential_learning}

A growing body of work explores how agents can accumulate and leverage experience over extended interactions, moving beyond static policy deployment toward systems that genuinely learn through doing.

AriGraph~\citep{anokhin2024arigraph} constructs knowledge graph world models with episodic memory for LLM agents, enabling them to maintain structured representations of their environment across interactions. This approach addresses a key limitation of standard LLM agents: the inability to retain and organize discoveries, spatial relationships, and causal knowledge beyond the current context window. Our PokeAgent harness adopts a similar philosophy through its persistent knowledge base, which accumulates facts, strategies, and observations across the agent's trajectory.

GRACE~\citep{song2024grace} introduces generalist agents that leverage both code generation and experiential learning, demonstrating that combining procedural knowledge (executable code) with episodic knowledge (accumulated experience) produces more robust agent behavior than either alone. This dual-knowledge paradigm aligns with AutoEvolve's design, which maintains both executable skills (Python code in a sandbox) and behavioral skills (textual strategies) alongside a curated memory store.

Language Models Meet World Models~\citep{xiang2023embodied} showed that grounding LLMs in embodied experience---rather than relying solely on pretraining knowledge---significantly improves their performance on downstream reasoning tasks. This finding motivates our emphasis on in-context learning from trajectory data: AutoEvolve's evolutionary optimization is driven entirely by the agent's own gameplay experience, using trajectory segments as the primary signal for harness refinement.

Recent work on in-context reinforcement learning (ICRL) has demonstrated that LLMs can perform implicit policy improvement from in-context reward signals~\citep{song2025icrl, karten2025economist}. These results provide theoretical grounding for our approach: if language models can implicitly improve policies from reward signals presented in-context, then explicitly structuring the agent's trajectory and reflection process should amplify this capability---which is precisely what AutoEvolve's evolutionary optimization loop achieves.

\subsubsection{Recursive Language Models}\label{sec:rlm}

Recursive Language Models (RLMs)~\citep{zhang2025rlm} formalize the paradigm of LLMs that programmatically interact with and recurse over their own context. Rather than treating the context window as a passive repository of information, RLMs treat it as an active workspace that the model can read from, write to, and restructure during inference. This provides a theoretical foundation for agents that self-modify their reasoning process---moving beyond simple chain-of-thought reasoning toward systems that can reorganize their own cognitive scaffolding.

The RLM framework is directly relevant to two components of our system. First, the \emph{trajectory window} in the PokeAgent harness (Chapter~4) functions as an RLM-style workspace: the full history of the agent's execution (state observations, prompts, reasoning traces, actions) is maintained as a structured, programmatically accessible log that subagents process and summarize for the orchestrator. This is a departure from standard agent architectures that treat history as a flat, append-only sequence. Second, AutoEvolve's mid-episode evolution (Chapter~5) can be understood as a form of recursive self-modification: the agent analyzes its own trajectory segments and uses this analysis to rewrite its harness components (system prompt, subagents, skills, memory), which in turn shape subsequent trajectory segments. Unlike the depth-$k$ recursion that RLMs formalize, AutoEvolve operates at depth-1 (the evolutionary optimizer analyzes the trajectory but does not recursively analyze its own optimization process), though the framework naturally extends to deeper recursion.

\subsubsection{Self-Adaptation}\label{sec:self_adaptation}

Several lines of work explore how language models and LLM-based agents can improve their own performance through self-directed learning, each making different tradeoffs between the scope of adaptation, the requirement for environment resets, and the computational overhead involved.

\paragraph{Prompt Evolution.} GEPA~\citep{agrawal2025gepa} uses genetic evolution over prompt candidates, guided by natural-language reflections on agent trajectories. It runs an agent multiple times, collects reasoning traces and outcomes, and has the model analyze these traces to propose prompt edits. Candidates that improve performance are retained for the next generation. GEPA demonstrated that this reflective evolution can outperform RL-based fine-tuning across several benchmarks, establishing prompt evolution as a viable alternative to weight updates. However, GEPA requires complete episodes and environment resets between iterations: each evolution step begins from the initial state, making it inapplicable to settings where resets are expensive or unavailable.

MIPROv2~\citep{opsahlong2024miprov2} applies multi-stage prompt optimization to DSPy programs, optimizing both instructions and demonstrations through Bayesian surrogate models. While effective for modular NLP pipelines, it assumes a fixed program structure and does not extend to the open-ended harness evolution that AutoEvolve targets.

\paragraph{Weight-Level Self-Adaptation.} Self-Adapting LLMs (SEAL)~\citep{zweiger2025seal} take an RL-first approach, enabling LLMs to generate their own fine-tuning data and apply policy-gradient updates to themselves. Over time, the model's weights evolve to incorporate new knowledge. While SEAL yields persistent improvements, it requires careful training to avoid catastrophic forgetting and operates at the weight level rather than the scaffold level. Our approach leaves model weights entirely unchanged and instead adapts via harness optimization---a strictly inference-time strategy that is simpler to deploy and leaves the base model available for other tasks.

\paragraph{Self-Refine.} The Self-Refine technique~\citep{madaan2023selfrefine} formalized the intuitive idea of iterative self-critique: the model produces an initial solution, generates feedback highlighting flaws, and revises its solution accordingly. This loop can be repeated to elicit higher-quality results without external feedback. Self-Refine operates within a single inference call and does not persist improvements across interactions. Our reflection and subagent mechanisms in PokeAgent (Chapter~4) extend this idea to the agent setting, where a dedicated reflecting subagent critiques the orchestrator's recent behavior and produces structured feedback that informs subsequent decisions.

\paragraph{Skill Learning.} Voyager's~\citep{wang2023voyager} skill library demonstrated that language model agents can progressively acquire reusable capabilities through interaction, storing successful programs in a growing code library indexed by natural language descriptions. GRACE~\citep{song2024grace} extended this with dual code-and-experience skills. AutoEvolve adopts a similar skill architecture but adds \emph{evolutionary curation}: skills are not only added but also updated, deprecated, or refined based on trajectory analysis during each evolution cycle.

\paragraph{Reset-Free Optimization.} A key distinction of our work is the \emph{reset-free} setting. Standard prompt evolution~\citep{agrawal2025gepa} and most RL-based methods require resetting the environment to a known initial state between optimization iterations. In the long-horizon RPG setting, this is problematic: a single speedrunning episode can last 12+ hours, and progress is inherently cumulative (items collected, battles won, areas unlocked). AutoEvolve evolves its harness mid-episode using trajectory \emph{segments} rather than complete episodes, enabling continuous adaptation without interrupting the agent's progress. This reset-free property distinguishes AutoEvolve from all prior prompt-evolution methods and is central to its applicability in embodied settings.

\paragraph{LLM-as-a-Judge.} The use of language models as automated evaluators has become widespread, with MT-Bench and Chatbot Arena~\citep{zheng2023mtbench} establishing the paradigm of LLM-as-a-Judge for open-ended evaluation. Subsequent work has identified systematic biases in this approach, including position bias and verbosity bias~\citep{zhao2024llmjudgefailures}, and has extended the paradigm to multimodal settings~\citep{chen2024mllmjudge}. AutoEvolve's evolutionary optimization relies on an LLM judge module that evaluates trajectory segments and produces structured critiques. We address known biases through careful prompt design and by grounding the judge's evaluations in concrete game-state metrics (milestone progress, step counts, spatial positioning) rather than relying solely on qualitative assessment.

\subsubsection{Meta-Harnesses}\label{sec:meta_harnesses}

The most directly relevant line of work treats the agent harness itself---the system prompts, tool definitions, retry logic, context management, and subagent configuration that wrap a foundation model---as an optimizable system. This emerging research direction recognizes that harness quality has become a first-order variable in agent performance, often mattering more than the choice of underlying model.

The Darwin G\"{o}del Machine (DGM)~\citep{zhang2025dgm} introduced open-ended evolution of self-improving coding agents. DGM maintains an archive of generated coding agents and grows the archive by sampling an agent and using a foundation model to create new, ``interesting'' variants---inspired by quality-diversity algorithms from open-endedness research. This process produces a growing tree of diverse, high-quality agents. Empirically, DGM improved coding capabilities on SWE-bench from 20.0\% to 50.0\% and on Polyglot from 14.2\% to 30.7\%, discovering improvements such as better code editing tools, long-context window management, and peer-review mechanisms. DGM operates across episodes (each agent variant is evaluated independently), requiring environment resets between iterations.

Meta-Harness~\citep{lee2026metaharness} automated harness optimization by giving a coding agent (Claude Code) access to a filesystem containing the full source code, execution traces, and performance scores of all prior harness candidates---up to 10 million tokens of diagnostic context per optimization step. This 400$\times$ increase over prior methods (which typically provided at most 26K tokens of summary information) enables the optimizer to trace failures to specific lines of harness code or prompt phrases. Meta-Harness achieved 76.4\% on Terminal-Bench 2.0 with Claude Opus 4.6 and demonstrated a 6$\times$ performance gap on the same underlying model through harness optimization alone.

AutoHarness~\citep{lou2026autoharness} demonstrated that Gemini-2.5-Flash can automatically synthesize code harnesses through iterative refinement guided by environment feedback, using tree search with Thompson sampling to explore the space of potential harnesses. The resulting harnesses prevent illegal moves in 145 TextArena games, enabling the smaller Gemini-2.5-Flash to outperform larger models including Gemini-2.5-Pro and GPT-5.2-High. AutoHarness focuses on \emph{constraint enforcement}---preventing illegal actions---rather than the broader harness evolution (prompts, subagents, skills, memory) that AutoEvolve targets.

Our AutoEvolve framework differs from all three approaches in several key respects. First, AutoEvolve is \emph{fully online and reset-free}: it evolves the harness during a single continuous episode rather than across independent evaluation runs. Second, it is \emph{tailored for embodied settings} with raw perceptual input (frame observations and button inputs), whereas DGM, Meta-Harness, and AutoHarness operate in textual environments with structured programmatic feedback. Third, AutoEvolve evolves the \emph{complete harness}---system prompt, subagent registry, skill library, and memory store---through four coordinated evolutionary passes, rather than optimizing a single component (code, prompts, or constraints) in isolation. These differences reflect the distinct demands of long-horizon embodied autonomy, where the agent cannot restart, perception is noisy, and the space of useful scaffolding is too large for any single-axis optimization to explore effectively.
